\chapter*{About this edition}
\addcontentsline{toc}{chapter}{About this edition}

The publication has three outputs from one modular LaTeX source tree: \textit{The Bijou Book}, \textit{Bijou Design Archive}, and \textit{Bijou Complete Edition}. The main book contains five teaching Parts, seven runnable workshops and six connecting chapters. Historical source records appear only in the archive and the complete edition's historical appendices. The main book includes a component atlas, renderer gallery, glossary and topic finder for practical reference.

All 426 source records and the frozen corpus survive this restructuring. Their prior typeset content, native TikZ figures and component-rendered tables remain unchanged. The 32 Mermaid conversions and 83 arrow-notation figures belong to those historical records; the two interaction traces belong to the workshops. Tables continue to use actual Bijou component output. Artwork comes from three explicitly distinct contexts: the 7.2.0 companion, the 8.0.0-rc.1 component capture checkout and the separate 7.1.0 renderer demonstrations.

Jacquard 24 supplies titles and all heading levels. IBM Plex Mono supplies prose, VT323 source listings, and Recursive Mono terminal specimens and component tables, with DejaVu Sans symbols. The existing coral palette is retained. Original font licenses, source hashes, capture parameters and code-example notices accompany the sources.

The version page records what was tested for the companion. New connecting chapters explain the same implementation and propose clearly labeled reader exercises. They do not claim that the example implements a process runner, a complete localization catalog, a built-in inspector or a verified human screen-reader workflow.
